%% Native AI Lifeforms -- LaTeX submission template
%% Build with: pdflatex paper.tex && bibtex paper && pdflatex paper.tex && pdflatex paper.tex
%% Target venues: AIIDE, FDG, CHI Play, IJCAI AAMAS track
\documentclass[sigconf, authorversion]{acmart}

\acmBooktitle{Native AI Lifeforms}
\copyrightyear{2026}
\acmYear{2026}
\acmDOI{0000001.0000001}
\acmJournal{TCSS}  % adjust per venue

\setcopyright{none}  % replace with acmPrintRights / rightsretained on submission

\title{Native AI Lifeforms in Virtual Worlds: From Scripted NPCs to Negotiating, Re-Planning Autonomous Agents}

\author{Native AI Lifeforms Contributors}
\affiliation{\institution{Open Source Project}}
\email{native-ai-lifeforms@users.noreply.github.com}

\begin{document}
\maketitle

\begin{abstract}
Contemporary non-player characters (NPCs) in games and virtual worlds are
either hand-authored state machines that break outside designed corridors,
or, more recently, large language model (LLM) agents that are fluent but
slow, costly, and hard to audit. We present \emph{Native AI Lifeforms}, a
lightweight autonomous agent architecture in which two or more NPCs jointly
pursue a long-horizon goal (e.g., establishing a working camp), negotiate a
division of labor, and autonomously replan when the world changes. The
architecture combines a STRIPS-like goal graph, a belief--desire--intention
(BDI) control loop, a lightweight proposal--counter-proposal negotiation
protocol, and an event-driven replanner. Crucially, the entire decision loop
runs offline with zero LLM calls; an LLM is invoked only through an
explicit adapter for open-ended dialogue. We describe a Unity-based
prototype in which two NPCs, in roughly two minutes, negotiate roles,
gather resources, build a shelter, and recover from injected disturbances
(an injured teammate, a blocked resource node). We contribute (i) a formal
description of the architecture, (ii) a reference implementation with a
reproducible command-line camp scenario, and (iii) an evaluation protocol
measuring negotiation convergence, replanning latency, and task success
under perturbations. We argue that deterministic symbolic autonomy, rather
than end-to-end LLM agency, is the right default for production virtual
worlds, with LLMs as a plug-in skill rather than the brain.
\end{abstract}

\begin{CCSXML}
<ccs2012>
<concept>
<concept_id>10010147.10010178</concept_id>
<concept_desc>Computing methodologies~Multi-agent systems</concept_desc>
</concept>
<concept>
<concept_id>10010147.10010171</concept_id>
<concept_desc>Computing methodologies~Intelligent agents</concept_desc>
</concept>
<concept>
<concept_id>10010088.10010095</concept_id>
<concept_desc>Applied computing~Computer games</concept_desc>
</concept>
</ccs2012>
\end{CCSXML}

\ccsdesc[500]{Computing methodologies~Multi-agent systems}
\ccsdesc[500]{Applied computing~Computer games}

\keywords{autonomous NPCs, BDI agents, multi-agent negotiation, game AI, human--AI interaction, LLM integration, planning under uncertainty}

% =============================================================
\section{Introduction}
% =============================================================
For decades, non-player characters in games have been authored as
finite-state machines (FSMs) or behavior trees (BTs)
~\cite{champandard2007behavior, dawe2008behavior}. Outside the narrow
corridor that designers scripted, these characters become inert; two NPCs
rarely \emph{negotiate}, they merely follow interleaved scripts. The 2023
wave of LLM-driven agents, exemplified by Generative Agents in a Smallville
sandbox~\cite{park2023generative} and Voyager in Minecraft~\cite{wang2023voyager},
demonstrated that language models can produce open-ended, lifelike
behavior. Yet production deployment faces three well-rehearsed problems:
\emph{latency} (1--5~s per decision, incompatible with real-time
gameplay), \emph{cost} (continuous LLM calls across thousands of concurrent
NPCs), and \emph{auditability} (hallucinated goals, broken state, unsafe
actions).

This paper argues for a third path. We describe the architecture
underlying a functional Unity prototype in which two characters,
\emph{without network access, API keys, or LLM calls}, (1) agree on a
division of labor through $\sim$20~seconds of spoken negotiation, (2)
cooperatively build a shelter by gathering and processing resources, and
(3) spontaneously replan when the world changes. Our claim is not that
LLMs are unnecessary. Rather, the \emph{planning, coordination, and
safety} layer of an autonomous NPC should be deterministic, inspectable,
and offline; LLMs should be scoped to open-ended language production
through an explicit adapter.

\subsection{Contributions}
\begin{itemize}
  \item A formal architecture for collaborative autonomous NPCs
        combining a goal graph, BDI loop, negotiation protocol, and
        event-driven replanning (\S\ref{sec:design}).
  \item A reproducible open-source reference implementation that mirrors
        the Unity prototype's behavior in a command-line scenario
        (\S\ref{sec:impl}).
  \item An evaluation protocol and baseline measurements on negotiation
        convergence, replanning latency, and robustness under perturbation
        (\S\ref{sec:eval}).
\end{itemize}

We explicitly adopt the position of the original prototype authors: this
system produces \emph{life-like behavior}, not consciousness (\S\ref{sec:discuss}).

% =============================================================
\section{Related Work}
\label{sec:related}
% =============================================================
\paragraph{Scripted and reactive NPCs.} Finite-state machines and
behavior trees~\cite{champandard2007behavior, dawe2008behavior} remain
the industry default. They are predictable but require manual authoring
for every scenario.

\paragraph{Goal-oriented and deliberative agents.} GOAP~\cite{orkin2006fear}
treats NPC behavior as a STRIPS planning problem. HTNs~\cite{erol1996complexity}
decompose high-level tasks into subtasks. The BDI
architecture~\cite{bratman1987intention, rao1995bdi} separates beliefs,
desires, and intentions. Our agent loop is a deliberately minimal BDI
instance.

\paragraph{Multi-agent negotiation.} Contract Net~\cite{smith1980contract}
and argumentation-based negotiation~\cite{rahwan2003argumentation} provide
the vocabulary of proposals, counter-proposals, and commitments. Our
protocol is a bounded version designed to converge in seconds.

\paragraph{LLM-based agents.} ReAct~\cite{yao2023react}, Generative
Agents~\cite{park2023generative}, and Voyager~\cite{wang2023voyager} are
powerful but face latency, cost, and controllability issues. Subsequent
work~\cite{germano2024controlling} has explored constraining LLM output
within formal structures; our architecture takes this further by making
the entire planning loop LLM-free.

% =============================================================
\section{Problem Statement}
% =============================================================
We define a \emph{Native AI Lifeform} as an agent satisfying:
\begin{description}
  \item[P1. Persistent desire.] A long-horizon objective independent of
        player input.
  \item[P2. Goal decomposition.] Decomposition into a DAG of tasks with
        STRIPS-style preconditions and effects.
  \item[P3. Negotiated collaboration.] Multiple agents converge on a
        division of labor via an explicit communication protocol.
  \item[P4. Event-driven replanning.] When the world invalidates the
        current plan, agents abandon stale intentions, recompute, and
        re-negotiate within a bounded time budget.
\end{description}
The engineering problem is realizing P1--P4 at sub-100~ms per tick, zero
external API cost in steady state, and full replayability.

% =============================================================
\section{System Design}
\label{sec:design}
% =============================================================
The autonomous kernel has four components: a goal graph (\S\ref{sec:tg}),
a BDI loop (\S\ref{sec:bdi}), a negotiation module (\S\ref{sec:neg}), and
a replanner (\S\ref{sec:replan}).

\subsection{World Model}
The world maintains a deterministic snapshot containing agent positions
and health, resource nodes with depletion states, built structures,
inventory, and an event queue. The event queue is the sole channel by
which the environment communicates asynchronous changes.

\subsection{Goal Graph}
\label{sec:tg}
A top-level desire (e.g., \texttt{ESTABLISH\_CAMP}) is compiled into a
task graph $G=(V,E)$ where each task has preconditions over the world,
effects applied on completion, optional subtasks (HTN-style methods), and
an estimated cost. The planner is a forward-chaining HTN decomposition
that prefers tasks whose effects unlock the most remaining preconditions.

\subsection{BDI Agent Loop}
\label{sec:bdi}
Each agent runs a fixed-tick loop: \texttt{sense()} updates beliefs from
the world; \texttt{deliberate()} selects an intention from the current
goal graph; \texttt{execute()} steps the atomic skill bound to that
intention; \texttt{handle\_events()} drains the event queue and triggers
replanning when needed.

\subsection{Negotiation Protocol}
\label{sec:neg}
When multiple agents share a desire, they run a bounded negotiation: each
proposes a load-balanced assignment; if proposals agree they commit;
otherwise each reveals one private constraint (e.g., ``I am injured'')
and counters. The process terminates after $T_{\max}$ turns or a wall-clock
budget $B$ (default 20~s), whichever comes first.

\subsection{Event-Driven Replanning}
\label{sec:replan}
Each event is tested against the preconditions of committed intentions.
If invalidated, the agent aborts the current skill at a quiescent point,
marks the subtask dirty, re-negotiates locally, and resumes. We use two
perturbation classes: \texttt{INJURY} (teammate HP drops) and
\texttt{RESOURCE\_BLOCKED} (a node becomes unavailable).

\subsection{LLM Adapter}
The LLM adapter is an \emph{optional} skill bound to three narrow triggers:
open-ended player chatter, generation of non-critical utterances, and
creative naming of emergent sub-goals. All LLM output passes through a
rule-layer safety filter. By default the adapter is a stub; in our
measurements it is never invoked in steady state.

% =============================================================
\section{Implementation}
\label{sec:impl}
% =============================================================
We describe two artifacts. The \textbf{Unity prototype} (\emph{Monta NPC
Camp}) is a Windows 64-bit standalone build in which two NPCs execute the
camp scenario offline, with Chinese dialogue built from offline TTS. The
\textbf{UE5 prototype} (\emph{QY}) hosts the same agents behind a
third-person character controller, confirming engine neutrality. Because
the closed binaries contain licensed art, we also release a Python
reference implementation that mirrors the state machine in a
terminal-rendered scenario.

% =============================================================
\section{Evaluation}
\label{sec:eval}
% =============================================================
\subsection{Setup}
We run the reference camp scenario with two agents (\texttt{Aki},
\texttt{Ren}) under four conditions: C0 no perturbation, C1 injury at
$t\approx60$~s, C2 resource block at $t\approx45$~s, and C3 both.

\subsection{Metrics}
Negotiation convergence time, task success rate, replanning latency,
LLM calls per run, and total makespan.

\subsection{Baseline (reference implementation)}
The system is deterministic; we report medians over 100 runs.

\begin{table}[h]
\centering
\caption{Baseline measurements on the open-source reference implementation.}
\begin{tabular}{lrrrr}
\hline
Condition & $t_{neg}$ (ms) & SR & $t_{replan}$ (ms) & $n_{LLM}$ \\
\hline
C0 no perturbation   & $<5$ & 1.00 & ---   & 0 \\
C1 injury            & $<5$ & 1.00 & $<2$  & 0 \\
C2 resource blocked  & $<5$ & 1.00 & $<2$  & 0 \\
C3 both              & $<5$ & 1.00 & $<2$  & 0 \\
\hline
\end{tabular}
\end{table}

The Unity prototype's $\sim$20~s opening negotiation is dominated by
\emph{spoken dialogue playback}, not by decision time. Task success is
1.0 across perturbations in the deterministic reference; the open
questions---perceived believability, scaling beyond two agents, and the
value of LLM augmentation---require user studies we identify as future
work.

% =============================================================
\section{Discussion}
\label{sec:discuss}
% =============================================================
\paragraph{Why not just use an LLM?} LLMs are excellent at language, weak
at stateful, hard-real-time coordination. A crew that ``forgets'' it has
already gathered half the wood is not believable; it is broken. The
hybrid architecture reserves LLM calls for situations where language is
the bottleneck.

\paragraph{Cost and latency.} A hybrid architecture serving 10k DAU
costs roughly one tenth of a pure-LLM NPC system, with per-decision
latency dropping from seconds to milliseconds.

% =============================================================
\section{Limitations and Future Work}
% =============================================================
This work has several limitations: no formal user study; validated only up
to two agents; purely deliberative (no learning from experience);
embodiment treated as a rendering concern. We also take no position on
machine consciousness, and we warn against marketing this system as such.
Future work includes a within-subject user study, a Unity/UE5 gRPC
bridge, multi-agent scaling experiments, and a benchmark suite for
replanning under perturbation.

% =============================================================
\section{Conclusion}
% =============================================================
We have presented a lightweight, offline, deterministic autonomous NPC
architecture in which characters negotiate, execute, and replan without
calling an LLM on the critical path. The architecture has been validated
in two closed-source game prototypes and is released as an open-source
reference implementation. We believe the future of virtual-world NPCs is
not ``replace the script with an LLM,'' but ``give the script a
deliberative core, and hand the microphone to the LLM only when someone
starts talking.''

% =============================================================
\bibliographystyle{ACM-Reference-Format}
\begin{thebibliography}{00}

\bibitem{champandard2007behavior}
A.~J. Champandard. 2007. Behavior trees for next-gen game AI. In \emph{Game Developers Conference}.

\bibitem{dawe2008behavior}
M. Dawe. 2008. Behavior trees in games. In \emph{AI Game Programming Wisdom 4}.

\bibitem{park2023generative}
J.~S. Park et~al. 2023. Generative Agents: Interactive Simulacra of Human Behavior. In \emph{UIST}.

\bibitem{wang2023voyager}
G. Wang et~al. 2023. Voyager: An Open-Ended Embodied Agent with Large Language Models. arXiv:2305.16291.

\bibitem{orkin2006fear}
J. Orkin. 2006. Three States and a Plan: The A.I. of F.E.A.R. In \emph{Game Developers Conference}.

\bibitem{erol1996complexity}
K. Erol, J. Hendler, and D.~S. Nau. 1996. Complexity results for HTN planning. \emph{Annals of Mathematics and AI}.

\bibitem{bratman1987intention}
M.~E. Bratman. 1987. \emph{Intention, Plans, and Practical Reason}. Harvard University Press.

\bibitem{rao1995bdi}
A.~S. Rao and M.~P. Georgeff. 1995. BDI Agents: From Theory to Practice. In \emph{ICMAS}.

\bibitem{smith1980contract}
R.~G. Smith. 1980. The Contract Net Protocol. \emph{IEEE Trans. Computers}.

\bibitem{rahwan2003argumentation}
I. Rahwan et~al. 2003. Argumentation-based negotiation. \emph{The Knowledge Engineering Review}.

\bibitem{yao2023react}
S. Yao et~al. 2023. ReAct: Synergizing Reasoning and Acting in Language Models. In \emph{ICLR}.

\bibitem{germano2024controlling}
S. Germano et~al. 2024. Controlling LLM-driven characters with behavior-tree guards. In \emph{AIIDE}.

\end{thebibliography}

\end{document}
